\chapter{Complexity — Book Model Basics}
%%% generated by the blueprint pipeline from TCSlib.Complexity.CircuitComplexity.BookModelBasics
\section{Overview}

Declarations of the Lean module \texttt{TCSlib.Complexity.CircuitComplexity.BookModelBasics}.

\section{Declarations}

\begin{definition}[Emission of a list of trees]
\lean{BoolCircuit.emitTrees}
\label{BoolCircuit.emitTrees}
\leanok
\uses{BoolCircuit.DAGGate, BoolCircuit.GateKind, BoolCircuit.emitTree}
Given a natural number $n$, a gate kind $k$, a list of lists of vertex indices, and an
existing list of gates, define the result by processing each vertex list in order with
the tree-emission procedure. The output is the accumulated gate list and, in the same
order, the vertex returned by that procedure for each input list.
\end{definition}

\begin{definition}[Tree Circuit Specification]
\lean{BoolCircuit.TreesSpec}
\label{BoolCircuit.TreesSpec}
\leanok
\uses{BoolCircuit.DAGGate, BoolCircuit.DAGGate.eval, BoolCircuit.GateKind, BoolCircuit.GatesAcyclic, BoolCircuit.runWith, BoolCircuit.treeCost, BoolCircuit.vertexValue}
Let $n$ be a natural number, let $k$ be an AND, OR, or NOT gate, let $vss$ be a finite
list of lists of natural-number vertex labels, let $gs$ be a finite list of Boolean
gates, and let $r=(r_1,r_2)$ consist of a finite gate list and a list of vertex labels.
Suppose $r_1$ is obtained by appending gates to $gs$ and is acyclic, $r_2$ has the same
length as $vss$, and every vertex label in $r_2$ is less than $n+|r_1|$. For every
Boolean assignment to the $n$ inputs, the values at the vertices in $r_2$, in order,
equal the evaluations on the initial circuit $gs$ of gates of kind $k$ whose respective
argument lists are the lists in $vss$. The total number of gates in $r_1$ is at most the
number in $gs$ plus the sum, over the lists $vs$ in $vss$, of the tree cost for $|vs|$
inputs; every gate appended after $gs$ has kind $k$, has pairwise distinct arguments,
and has at most two arguments.
\end{definition}

\begin{theorem}[Correctness and size bound for tree emission]
\lean{BoolCircuit.emitTrees_spec}
\label{BoolCircuit.emitTrees_spec}
\leanok
\uses{BoolCircuit.DAGGate, BoolCircuit.DAGGate.eval_congr, BoolCircuit.GateKind, BoolCircuit.GatesAcyclic, BoolCircuit.TreesSpec, BoolCircuit.emitTree, BoolCircuit.emitTree_spec, BoolCircuit.emitTrees, BoolCircuit.runWith_append, BoolCircuit.runWith_getD_of_lt, BoolCircuit.vertexValue_append}
\difficulty{6}
Let $n$ be a natural number, let $k$ be an AND or OR gate kind, and let $V_1,\ldots,V_r$
be lists of vertex indices. Let $G$ be an acyclic list of Boolean gates, and suppose
every index in every $V_j$ is less than $n+|G|$. Emitting a tree of gates of kind $k$
for each list returns an extended acyclic gate list and one output index per list, each
less than $n$ plus the length of the resulting gate list. On every Boolean input to the
original $n$ vertices, the output values equal the values of the corresponding single
gates of kind $k$ applied to the vertices in each original list. The resulting gate list
has length at most $|G|+\sum_{j=1}^r c(|V_j|)$, where $c(0)=1$ and $c(s)=s-1$ for $s>0$;
every appended gate has kind $k$, has distinct input indices, and reads at most two
vertices.
\end{theorem}

\begin{definition}[Conjunction circuit]
\lean{BoolCircuit.andAllDAG}
\label{BoolCircuit.andAllDAG}
\leanok
\uses{BoolCircuit.DAGCircuit, BoolCircuit.GatesAcyclic.nil, BoolCircuit.emitTree, BoolCircuit.emitTree_spec}
For each natural number $n$, let $C_n$ be the directed acyclic Boolean circuit obtained
by emitting an AND tree on the input vertices $0,1,0,dots,n-1$; its output is the
vertex returned by that tree construction.
\end{definition}

\begin{theorem}[Specification of the emitted AND circuit]
\lean{BoolCircuit.andAllDAG_spec}
\label{BoolCircuit.andAllDAG_spec}
\leanok
\uses{BoolCircuit.GatesAcyclic.nil, BoolCircuit.TreeSpec, BoolCircuit.emitTree, BoolCircuit.emitTree_spec}
\difficulty{1}
For every natural number $n$, consider the circuit obtained by recursively combining the
input vertices $0,1,dots,n-1$ with AND, starting from an empty gate list, and let $r$
be its gate list and output vertex. The gate list extends the empty list and is acyclic,
and the output vertex is less than $n$ plus the number of gates; on every Boolean input
assignment, the output equals the AND of all input values. The number of gates is at
most the tree cost of $n$, and the output depth is at most $\lceil\log_2 n\rceil+1$.
Every gate in the emitted list has AND kind, has distinct input vertices, and has at
most two inputs.
\end{theorem}

\begin{theorem}[Evaluation of the AND circuit]
\lean{BoolCircuit.andAllDAG_eval}
\label{BoolCircuit.andAllDAG_eval}
\leanok
\uses{BoolCircuit.DAGCircuit.eval, BoolCircuit.DAGCircuit.values, BoolCircuit.DAGGate.eval, BoolCircuit.andAllDAG, BoolCircuit.andAllDAG_spec, BoolCircuit.runWith_nil, BoolCircuit.vertexValue}
\difficulty{4}
For every natural number $n$ and every Boolean input assignment $x$ on the $n$ input
vertices, the circuit formed by combining all inputs with AND evaluates to true if and
only if every input is true.
\end{theorem}

\begin{theorem}[Fan-in bound for the AND circuit]
\lean{BoolCircuit.andAllDAG_isFaninTwo}
\label{BoolCircuit.andAllDAG_isFaninTwo}
\leanok
\uses{BoolCircuit.DAGCircuit.IsFaninTwo, BoolCircuit.andAllDAG, BoolCircuit.andAllDAG_spec}
\difficulty{3}
For every natural number $n$, the circuit formed by combining all $n$ inputs with AND is
well formed: each gate has distinct input vertices, and every NOT gate has exactly one
input. Moreover, every gate has at most two input vertices.
\end{theorem}

\begin{theorem}[Size bound for the AND circuit]
\lean{BoolCircuit.andAllDAG_size_le}
\label{BoolCircuit.andAllDAG_size_le}
\leanok
\uses{BoolCircuit.DAGCircuit.size, BoolCircuit.andAllDAG, BoolCircuit.andAllDAG_spec, BoolCircuit.emitTree, BoolCircuit.treeCost, BoolCircuit.treeCost_le}
\difficulty{3}
For every natural number $n$, the size of the circuit formed by combining all $n$ inputs
with AND is at most $2n+1$.
\end{theorem}

\begin{definition}[Family of conjunction circuits]
\lean{BoolCircuit.andAllFamily}
\label{BoolCircuit.andAllFamily}
\leanok
\uses{BoolCircuit.DAGCircuitFamily, BoolCircuit.andAllDAG}
There is a family of directed acyclic Boolean circuits that, for each natural number
$n$, assigns the conjunction circuit obtained by emitting an AND tree on the $n$ input
vertices.
\end{definition}

\begin{theorem}[Linear-size circuits for the all-ones language]
\lean{Language.allOnes_inSIZE_linear}
\label{Language.allOnes_inSIZE_linear}
\leanok
\uses{BoolCircuit.DAGCircuit.eval, BoolCircuit.DAGCircuitFamily.mem_language_iff, BoolCircuit.andAllDAG, BoolCircuit.andAllDAG_eval, BoolCircuit.andAllDAG_isFaninTwo, BoolCircuit.andAllDAG_size_le, BoolCircuit.andAllFamily, Language.InSIZE, Language.allOnes}
\difficulty{4}
The language consisting of all finite Boolean strings whose entries are all true is
recognized by a family of Boolean circuits in which every gate is well formed and has at
most two inputs, and the circuit for inputs of length $n$ has size at most $2n+1$.
\end{theorem}

\begin{theorem}[Polynomial size for the all-ones language]
\lean{Language.allOnes_inPPoly}
\label{Language.allOnes_inPPoly}
\leanok
\uses{Language.InPPoly, Language.InSIZE.inPPoly, Language.allOnes, Language.allOnes_inSIZE_linear}
\difficulty{3}
The language consisting of Boolean words whose entries are all true is recognized by a
family of well-formed acyclic Boolean circuits with fan-in at most two, with size
bounded by a polynomial in the input length.
\end{theorem}

\begin{theorem}[Tree Cost Bound]
\lean{BoolCircuit.treeCost_le_max}
\label{BoolCircuit.treeCost_le_max}
\leanok
\uses{BoolCircuit.treeCost}
\difficulty{3}
For natural numbers $m$ and $k$ with $m\le k$, the tree cost of $m$ is at most
$\max(1,k-1)$.
\end{theorem}

\begin{theorem}[Length bound for gate binarization]
\lean{BoolCircuit.rewriteGates_binarize_length_le}
\label{BoolCircuit.rewriteGates_binarize_length_le}
\leanok
\uses{BoolCircuit.DAGGate, BoolCircuit.GatesAcyclic, BoolCircuit.binarizeGadget, BoolCircuit.binarizeGadget_correct, BoolCircuit.emitTree_spec, BoolCircuit.rewriteGates, BoolCircuit.rewriteGates_inv, BoolCircuit.rewriteStep, BoolCircuit.treeCost_le_max}
\difficulty{6}
Fix a natural number $K$ and an acyclic list of gates on $n$ input vertices. Suppose
every gate has at most $K$ arguments and every NOT gate has exactly one argument.
Rewriting the list by replacing each gate with its binarized gadget produces a list
whose length is at most the sum, over the original gates, of the larger of $1$ and the
number of arguments minus $1$.
\end{theorem}

\begin{theorem}[Refined size bound for binarization]
\lean{BoolCircuit.DAGCircuit.binarize_size_le_sum}
\label{BoolCircuit.DAGCircuit.binarize_size_le_sum}
\leanok
\uses{BoolCircuit.DAGCircuit, BoolCircuit.DAGCircuit.IsWellFormed, BoolCircuit.DAGCircuit.args_length_le_size, BoolCircuit.DAGCircuit.binarize, BoolCircuit.DAGCircuit.size, BoolCircuit.binarizeGadget, BoolCircuit.rewriteGates, BoolCircuit.rewriteGates_binarize_length_le}
\difficulty{3}
Let $C$ be a directed acyclic Boolean circuit with $n$ input vertices, and suppose each
gate has pairwise distinct arguments and every negation gate has exactly one argument.
The size of the binarized circuit is at most $n+\sum_{g\in G}\max\{1,\ell(g)-1\}$, where
$G$ is the gate list and $\ell(g)$ is the number of arguments of gate $g$.
\end{theorem}

\begin{theorem}[Binarization with a size bound]
\lean{BoolCircuit.DAGCircuit.exists_faninTwo_size_le_sum}
\label{BoolCircuit.DAGCircuit.exists_faninTwo_size_le_sum}
\leanok
\uses{BoolCircuit.DAGCircuit, BoolCircuit.DAGCircuit.IsFaninTwo, BoolCircuit.DAGCircuit.IsWellFormed, BoolCircuit.DAGCircuit.binarize, BoolCircuit.DAGCircuit.binarize_eval, BoolCircuit.DAGCircuit.binarize_isFaninTwo, BoolCircuit.DAGCircuit.binarize_size_le_sum, BoolCircuit.DAGCircuit.eval, BoolCircuit.DAGCircuit.size}
\difficulty{1}
Let $C$ be a directed acyclic Boolean circuit on $n$ inputs such that each gate has no
repeated arguments and every NOT gate has exactly one argument. There is a circuit $C'$
on the same inputs that satisfies these same conditions and has at most two arguments at
every gate, computes the same Boolean function as $C$, and has size at most
$n+\sum_{g\in C}\max(1,\ell(g)-1)$, where $\ell(g)$ is the number of arguments of $g$.
\end{theorem}

\begin{definition}[Input negation gates]
\lean{BoolCircuit.negInputGates}
\label{BoolCircuit.negInputGates}
\leanok
\uses{BoolCircuit.DAGGate}
For each natural number $n$, form the list of $n$ NOT gates, with the gate at position
$k$ reading input vertex $k$, for $0 \le k < n$.
\end{definition}

\begin{theorem}[Number of negation input gates]
\lean{BoolCircuit.length_negInputGates}
\label{BoolCircuit.length_negInputGates}
\leanok
\uses{BoolCircuit.negInputGates}
\difficulty{1}
For every $n\in\mathbb{N}$, the list of negation gates on the $n$ input vertices has
length $n$.
\end{theorem}

\begin{theorem}[Acyclicity of input-negation gates]
\lean{BoolCircuit.negInputGates_acyclic}
\label{BoolCircuit.negInputGates_acyclic}
\leanok
\uses{BoolCircuit.GatesAcyclic, BoolCircuit.length_negInputGates, BoolCircuit.negInputGates}
\difficulty{3}
For every nonnegative integer $n$, the list of $n$ input-negation gates is acyclic: at
position $i$ in the list, every input vertex has index strictly less than $n+i$.
\end{theorem}

\begin{theorem}[Values of negated inputs]
\lean{BoolCircuit.vertexValue_negInputGates}
\label{BoolCircuit.vertexValue_negInputGates}
\leanok
\uses{BoolCircuit.DAGGate.eval, BoolCircuit.negInputGates, BoolCircuit.runWith_getD_gate, BoolCircuit.runWith_getD_of_lt, BoolCircuit.vertexValue}
\difficulty{4}
Let $n$ be a natural number, let $x$ assign Boolean values to the $n$ inputs, and let
$i$ be an input index. In the circuit consisting of one negation gate for each input,
the value at vertex $n+i$ is $
eg x(i)$.
\end{theorem}

\begin{definition}[Literal vertex encoding]
\lean{BoolCircuit.litVertex}
\label{BoolCircuit.litVertex}
\leanok
For a natural number $n$, an index $i$ among $0,\ldots,n-1$, and a Boolean value $b$,
assign the literal vertex $i$ when $b$ is true and the vertex $n+i$ when $b$ is false.
\end{definition}

\begin{theorem}[Literal vertex bound]
\lean{BoolCircuit.litVertex_lt}
\label{BoolCircuit.litVertex_lt}
\leanok
\uses{BoolCircuit.litVertex}
\difficulty{4}
Let $n\in\mathbb{N}$, let $i$ be an element of $\mathrm{Fin}(n)$, and let $b$ be a
Boolean value. The vertex assigned to the literal with index $i$ and polarity $b$ is
less than $2n$.
\end{theorem}

\begin{theorem}[Value of a literal vertex]
\lean{BoolCircuit.vertexValue_litVertex}
\label{BoolCircuit.vertexValue_litVertex}
\leanok
\uses{BoolCircuit.litVertex, BoolCircuit.negInputGates, BoolCircuit.vertexValue, BoolCircuit.vertexValue_input, BoolCircuit.vertexValue_negInputGates}
\difficulty{4}
Let $n$ be a natural number, let $x$ assign Boolean values to the $n$ inputs, let $i$ be
an input index, and let $b$ be a Boolean value. In the circuit consisting of one
negation gate for each input, the value at the vertex representing the literal that is
true exactly when $x(i)=b$ is the truth value of $x(i)=b$.
\end{theorem}

\begin{definition}[Minterm input vertices]
\lean{BoolCircuit.mintermLeaves}
\label{BoolCircuit.mintermLeaves}
\leanok
\uses{BoolCircuit.litVertex}
Given an assignment $v$ of Boolean values to $n$ inputs, form the list whose entry for
each input $i$ is its literal vertex: $i$ if $v(i)$ is true, and $n+i$ if $v(i)$ is
false.
\end{definition}

\begin{theorem}[Evaluation of minterm leaves]
\lean{BoolCircuit.mintermLeaves_eval}
\label{BoolCircuit.mintermLeaves_eval}
\leanok
\uses{BoolCircuit.DAGGate, BoolCircuit.DAGGate.eval, BoolCircuit.litVertex, BoolCircuit.mintermLeaves, BoolCircuit.negInputGates, BoolCircuit.runWith, BoolCircuit.vertexValue, BoolCircuit.vertexValue_litVertex}
\difficulty{4}
Let $n$ be a nonnegative integer and let $v,x : \mathrm{Fin}(n) \to
\{\mathrm{false},\mathrm{true}\}$. The AND gate whose inputs are the minterm leaves
specified by $v$, evaluated after the input-negation gates have been run on $x$, has
value equal to the Boolean indicator of the equality $x=v$.
\end{theorem}

\begin{definition}[Satisfying assignments list]
\lean{BoolCircuit.satList}
\label{BoolCircuit.satList}
\leanok
For a Boolean function $f$ on assignments to $n$ inputs, form the list of all
assignments $v$ for which $f(v)$ is true.
\end{definition}

\begin{definition}[Gate construction for disjunctive normal form]
\lean{BoolCircuit.dnfGates}
\label{BoolCircuit.dnfGates}
\leanok
\uses{BoolCircuit.DAGGate, BoolCircuit.emitTree, BoolCircuit.emitTrees, BoolCircuit.mintermLeaves, BoolCircuit.negInputGates, BoolCircuit.satList}
For a Boolean function $f\colon
(\mathrm{Fin}\,n\to\{\mathrm{false},\mathrm{true}\})\to\{\mathrm{false},\mathrm{true}\}$,
define its gate output to be the pair obtained by emitting AND trees from the minterm
leaves of the satisfying assignments of $f$, initialized with the input negation gates,
and then emitting an OR tree from the resulting roots. The pair consists of the
resulting gate list and its output vertex.
\end{definition}

\begin{theorem}[Specification of emitted minterms]
\lean{BoolCircuit.minterms_spec}
\label{BoolCircuit.minterms_spec}
\leanok
\uses{BoolCircuit.TreesSpec, BoolCircuit.emitTrees, BoolCircuit.emitTrees_spec, BoolCircuit.litVertex_lt, BoolCircuit.mintermLeaves, BoolCircuit.negInputGates, BoolCircuit.negInputGates_acyclic, BoolCircuit.satList}
\difficulty{4}
Let $n$ be a nonnegative integer and let $f$ be any Boolean function on assignments to
$n$ variables. Run the input-negation gates first, then emit one AND tree for each
assignment on which $f$ is true, using as that tree's leaves the literals selected by
the assignment. The resulting gates extend the input-negation gates and are acyclic; the
output list has one vertex per satisfying assignment, each output vertex is less than
$n$ plus the total number of resulting gates, and the value at each output equals the
corresponding AND of minterm leaves evaluated after the input-negation gates. The total
number of gates is at most $n$ plus the sum, over the satisfying assignments, of the
tree cost of a list of $n$ leaves; every newly emitted gate is an AND gate with distinct
inputs and at most two inputs.
\end{theorem}

\begin{theorem}[Specification of the DNF construction]
\lean{BoolCircuit.dnf_spec}
\label{BoolCircuit.dnf_spec}
\leanok
\uses{BoolCircuit.TreeSpec, BoolCircuit.dnfGates, BoolCircuit.emitTree_spec, BoolCircuit.emitTrees, BoolCircuit.mintermLeaves, BoolCircuit.minterms_spec, BoolCircuit.negInputGates, BoolCircuit.satList}
\difficulty{1}
For every natural number $n$ and Boolean function $f\colon\{0,1\}^n\to\{0,1\}$, the DNF
construction satisfies the following specification. Starting with one negation gate for
each input, it first emits conjunction trees for the minterm leaves of each satisfying
assignment, then emits an OR tree on the resulting vertices; the final gate list extends
the initial negation gates and is acyclic, and its output is a valid vertex. On every
input, that output has the same value as an OR gate applied to the values of the
satisfying-assignment minterms, evaluated using the initial negation gates. The final
list has at most the initial list length plus the tree cost of the number of those
minterms, and the output depth is at most the ceiling of the base-two logarithm of that
number plus one plus the maximum depth of the minterm vertices in the initial list.
Every gate added after the initial list is an OR gate with distinct arguments and at
most two inputs.
\end{theorem}

\begin{definition}[Disjunctive normal form circuit]
\lean{BoolCircuit.dnfDAG}
\label{BoolCircuit.dnfDAG}
\leanok
\uses{BoolCircuit.DAGCircuit, BoolCircuit.dnfGates, BoolCircuit.dnf_spec}
For a Boolean function $f$ on assignments to $n$ inputs, form the directed acyclic
circuit whose gates and output vertex are those obtained by constructing the disjunctive
normal form gate list and output.
\end{definition}

\begin{theorem}[Evaluation of the DNF circuit]
\lean{BoolCircuit.dnfDAG_eval}
\label{BoolCircuit.dnfDAG_eval}
\leanok
\uses{BoolCircuit.DAGCircuit.eval, BoolCircuit.DAGGate, BoolCircuit.DAGGate.eval, BoolCircuit.dnfDAG, BoolCircuit.dnfGates, BoolCircuit.dnf_spec, BoolCircuit.emitTrees, BoolCircuit.mintermLeaves, BoolCircuit.mintermLeaves_eval, BoolCircuit.minterms_spec, BoolCircuit.negInputGates, BoolCircuit.runWith, BoolCircuit.satList, BoolCircuit.vertexValue}
\difficulty{4}
Let $n$ be a natural number, let $f$ be a Boolean function on assignments to $n$ Boolean
variables, and let $x$ be such an assignment. The circuit constructed from the
satisfying assignments of $f$ evaluates at $x$ to $f(x)$.
\end{theorem}

\begin{theorem}[Fan-in of the DNF circuit]
\lean{BoolCircuit.dnfDAG_isFaninTwo}
\label{BoolCircuit.dnfDAG_isFaninTwo}
\leanok
\uses{BoolCircuit.DAGCircuit.IsFaninTwo, BoolCircuit.dnfDAG, BoolCircuit.dnfGates, BoolCircuit.dnf_spec, BoolCircuit.minterms_spec}
\difficulty{5}
For every natural number $n$ and every Boolean function $f\colon\{0,1\}^n\to\{0,1\}$,
the circuit obtained by the specified DNF construction is well formed: each gate reads
distinct vertices, and every negation gate has exactly one input. Every gate in the
circuit has at most two inputs.
\end{theorem}

\begin{theorem}[Bound on tree construction cost]
\lean{BoolCircuit.dnf_cost_le}
\label{BoolCircuit.dnf_cost_le}
\leanok
\uses{BoolCircuit.treeCost}
\difficulty{4}
Let $n$ and $s$ be natural numbers with $s \le 2^n$. Then $s$ times the tree cost for
$n$, plus the tree cost for $s$, is at most $n2^n+1$.
\end{theorem}

\begin{theorem}[Size of the DNF circuit]
\lean{BoolCircuit.dnfDAG_size_le}
\label{BoolCircuit.dnfDAG_size_le}
\leanok
\uses{BoolCircuit.DAGCircuit.size, BoolCircuit.dnfDAG, BoolCircuit.dnfGates, BoolCircuit.dnf_cost_le, BoolCircuit.dnf_spec, BoolCircuit.length_negInputGates, BoolCircuit.mintermLeaves, BoolCircuit.minterms_spec, BoolCircuit.satList, BoolCircuit.treeCost}
\difficulty{4}
For every natural number $n$ and every Boolean function $f\colon\{0,1\}^n\to\{0,1\}$,
the size of the DNF circuit constructed for $f$ is at most $n2^n+2n+1$.
\end{theorem}

\begin{theorem}[Size bound for a fan-in-two circuit]
\lean{BoolCircuit.exists_dagCircuit_faninTwo_size_le}
\label{BoolCircuit.exists_dagCircuit_faninTwo_size_le}
\leanok
\uses{BoolCircuit.DAGCircuit, BoolCircuit.DAGCircuit.IsFaninTwo, BoolCircuit.DAGCircuit.eval, BoolCircuit.DAGCircuit.size, BoolCircuit.dnfDAG, BoolCircuit.dnfDAG_eval, BoolCircuit.dnfDAG_isFaninTwo, BoolCircuit.dnfDAG_size_le}
\difficulty{1}
For every nonnegative integer $n$ and every Boolean function on assignments to $n$
variables, there is a directed acyclic circuit with those inputs and a valid output
vertex that computes the function on every assignment. Each gate has at most two inputs,
has no repeated input index, and every NOT gate has exactly one input; the circuit size,
counting the input vertices and gates, is at most $n2^n+2n+1$.
\end{theorem}
